% !TeX root = ../教材题库_全部条目_转写工作稿.tex
% Source: XJ-B1-C05-PR-P161-Y0388
% 第5章 / 5.1 任意角和弧度制 / 教材 p.161 / PDF 第168页
% OCR未校对：公式与文字均可能有误；未完成数学TeX转写；禁止直接发给学生
% 原题及相邻原解截图（核对用）：../题目原图/XJ-B1-C05-PR-P161-Y0388-01.webp
% 本文件只在转写完成后作为题干使用；OCR 条目尚不是可靠题干。
\begin{bankitem}{p.161 练习 第2题}{第5章 / 5.1 任意角和弧度制 / 教材 p.161 / PDF 第168页}{XJ-B1-C05-PR-P161-Y0388}{OCR未校对：公式与文字均可能有误；未完成数学TeX转写；禁止直接发给学生}
\begin{ocrdraft}
[教材p.161 / PDF第168页]
\par
2．用弧度表示：
\par
（1)终边在x轴上的角的集合；
\par
（2）终边在3，轴上的角的集合．
\end{ocrdraft}
\end{bankitem}
